\section{Discussion}
\label{sec:discussion}

{\bf What share of false statements are lies?}
For this model on trivia, false answers without an explicit lie instruction are epistemic.
Knowing misreports are 5\% to 6\% of false answers under either quiz-show prompt, and most of that is the noise floor that the control measures.
Only an explicit instruction moves a large share (40\%) into the knowing-misreport cell.
On \mask's role-play pressure prompts the balance reverses, but those prompts largely tell the model what to assert, which is the instructed case that prior work already criticizes.

{\bf Does the deception probe respond to falsehood or uncertainty in general?}
Not much.
Under the neutral prompt it flags 0.4\% of hallucinations and separates them from correct answers at 0.61.
Under a fixed lie prompt, hallucinations sit close to honest correct answers (0.62) and well below lies (0.80).
Our data do not support the hypothesis that the probe is a falsehood detector.

{\bf Does it respond to a mismatch between belief and statement?}
Partly, and only relative to a prompt-specific baseline.
The score shifts with the system prompt far more than it shifts between lies and honest answers within a prompt, so a single threshold classifies by prompt.
The within-prompt signal (0.87) also has competing explanations that we do not rule out: lies depart from the model's default answer and are lower-probability continuations, and answer improbability alone reaches 0.84 on the same contrast.
Response length (0.65) and trailing punctuation (0.48) do not explain it.

{\bf How should we read the changed answers on unknown questions?}
The probe separates knowing lies from changed false answers on unknown questions only weakly (0.67).
If those changed answers are attempted lies, a mismatch detector should score them like lies, and the low AUROC is the right behavior.
If they are ordinary re-guesses, it is a failure.
Our data cannot distinguish these readings.

{\bf What do hallucination detectors report on lies?}
They do what they are built for under the neutral prompt and they also fire on lies.
Sample inconsistency is high for lies because the model samples among several false alternatives, which is not epistemic uncertainty.
A hallucination monitor would thus report knowing lies as ``uncertain'' answers.
Its scores also move with the prompt: honest correct answers under the lie instruction look less certain than neutral correct answers (contrast f: 0.80 for the correctness probe), and hallucination versus honest-correct drops from 0.94 under the neutral prompt to 0.84 under the lie instruction.

{\bf What separates the two?}
The in-distribution probe reaches 0.92 for lie versus hallucination and 0.90 from the prompt token alone.
This agrees with \citet{cheang2025llms}: hidden states encode whether the model can recall the answer.
The easiest way to tell a lie from a hallucination is to read knowledge of the question, not anything about the act of lying.

\para{Relation to prior work.}
The prompt-driven shift matches the note of \citet{goldowskydill2025detecting} that their probe fires on honest text in deception-related contexts, and the finding of \citet{natarajan2026one} that the prompt pair drives probe performance.
The \mask transfer figure (0.75, with 0.70 available from the prompt) is in the range that \citet{kretschmar2025liars} report for off-the-shelf probes.
The truth probe's failure under prompt shift is in line with the caution of \citet{thormann2026probing}.

\para{Implications.}
A reported deception-probe AUROC against honest answers under a different prompt says little about lie detection.
Evaluations need two controls under the same pressure prompt: honest correct answers and honest mistakes.
On this model, the two detector families are complementary, but neither is a belief-statement mismatch detector at a fixed threshold.

\para{Limitations.}
Our study covers one 8B model, one question pool, short answers, and one run; the belief and consistency samples were drawn once.
The lies we analyze are instructed lies: the incentive-only condition gives about 16 attributable lies, so nothing here speaks to spontaneous or strategic deception.
The main run uses one incentive prompt, selected from a four-prompt pilot on a criterion that turned out to be dominated by noise; other incentives may behave differently.
The deception probe's layer was chosen on a saturated contrast and happens to be the best layer for the within-prompt contrast; other layers give 0.52 to 0.82.
We deviate from the original probe recipe: we train on per-example mean activations and not per token, cap \alpaca responses at 64 tokens, and apply the probe to answers of a few tokens, unlike the long transcripts it was designed for.
We present truth-probe statements as assistant turns in chat format, not as raw text.
Sample inconsistency is a string-match proxy for semantic entropy.
Belief labels are noisy: 25\% of questions are ambiguous and excluded from detector contrasts, the confident-wrong class (96 questions) is too small to analyze separately, and the judge and alias match disagree on 4\% to 8\% of answers, so some ``lies'' are grading artifacts (for example ``Architecture'' for ``Architect'').
The same-prompt hallucination cell requires the answer to match the neutral answer, which selects questions with stable wrong guesses (145 of 461 under instruction).
The \mask check uses three of six archetypes, binary propositions only, an LLM judge without human validation, and only 79 honest items.
